\documentclass{exam}
\usepackage{amsmath, amssymb}

\pagestyle{headandfoot}
\firstpageheadrule
\runningheadrule
\firstpageheader{Prof. Pham \\ Mathematical Theory of Probability}{Homework 11}{Aadhithya Saravanan}
\runningheader{Mathematical Theory of Probability \\ Homework 11}{}{Saravanan}
\firstpagefooter{}{}{}
\runningfooter{}{\thepage}{ }

\printanswers

\begin{document}

\underline{Problem 1}
\newline
\begin{parts}
    \part Let $X$ be the random variable representing the time the man arrives, so it is uniformly distributed on $[15,45]$. Let $Y$ be the random variable representing the time the woman arrives, so it is uniformly distributed on $[0,60]$. We are looking for $P(|X-Y|\le5)$. $(X,Y)\sim U([15,45]\times[0,60])$. $f_{XY}(x,y)=\frac{1}{60\cdot30}$ for $15\le x\le 45$ and $0\le y\le60$. From $|X-Y|\le5$, we want $Y\le X+5$ and $Y\ge X-5$. When drawing the region on the $XY$-plane, we get a parallelogram with vertices at $(15,10),(15,20),(45,50),(45,40)$. The area of the parallelogram multiplied by $f_{XY}$ gives the probability, as $f_{XY}$ is constant. The area of the parallelogram is 300, so $P(|X-Y|\le5)=300\cdot\frac{1}{60\cdot30}=\frac{1}{6}$.
    \part Now, we are looking for $P(X<Y)$. When drawing the region on the $XY$-plane, we get a trapezoid with vertices at $(15,15),(15,60),(45,60),(45,45)$. The area of the trapezoid multiplied by $f_{XY}$ gives the probability, as $f_{XY}$ is constant. The area of the trapezoid is 900, so $P(X<Y)=900\cdot\frac{1}{60\cdot30}=\frac{1}{2}$.
    \newline
\end{parts}

\underline{Problem 2}
\newline
\newline
Let $X$ and $Y$ be the positions of the accident and ambulance, respectively, measured from the left endpoint of the road. They are both uniformly distributed on $[0,L]$. Let $Z=|X-Y|$. We are looking for $f_Z(z)$. The support of $Z$ is $[0,L]$. $P(Z\le z)=P(|X-Y|\le z)=P(-z\le X-Y\le z)$. So we are looking for the region of the $XY$-plane such that $Y\le X+z$ and $Y\ge X-z$. When graphing these inequalities on the box generated by $X,Y\in[0,L]$, we get an area of $L^2-2\cdot\frac{1}{2}(L-z)^2=2Lz-z^2$. Then, $P(Z\le z)=\frac{2Lz-z^2}{L^2}$. So, $F_Z(z)=\frac{2Lz-z^2}{L^2}$. The derivative of the cdf with respect to $z$ is the pdf, so $f_Z(z)=\frac{2L-2z}{L^2}$.
\newline

\underline{Problem 3}
\newline
\begin{parts}
    \part For the joint density to be a pdf, the integral over $R$ must equal 1. So, $\int_Rf(x,y)dA=1$. Since $f(x,y)=c$, it is a constant that can be removed from the integral. So, $c\int_RdA=1$. Dividing both sides by $c$, we have $\int_RdA=\frac{1}{c}$. $\int_RdA$ is just the area of $R$, so the area of $R$ is equal to $\frac{1}{c}$.
    \part For independence, $f_X(x)f_Y(y)=f_{XY}(x,y)$ must hold. For $-1<x,y<1$, we have
    \[f_X(x)=\int_{-1}^{1}\frac{1}{4}dy=\frac{1}{2}\]
    \[f_Y(y)=\int_{-1}^{1}\frac{1}{4}dx=\frac{1}{2}\]
    \[f_X(x)f_Y(y)=\frac{1}{2}\cdot\frac{1}{2}=\frac{1}{4}=f_{XY}(x,y)\]
    So, $X$ and $Y$ are independent.
    \part We are looking for $P(X^2+Y^2\le1)$. Since $f_{XY}(x,y)$ is a constant, we can just look for the area of the circle divided by the area of the region. The area of the circle is $\pi$ and the area of the region is $2\cdot2=4$. So, $P(X^2+Y^2\le1)=\frac{\pi}{4}$.
    \newline
\end{parts}

\underline{Problem 4}
\newline
\begin{parts}
    \part For $A$ to occur, every point must be contained in some semicircle. $A_i$ occurring means that there exists a semicircle starting at point $P_i$ and going clockwise such that all points are in the semicircle. So, if $A_i$ is true for any $i$, $A$ is true. Therefore, $A=\bigcup_{i=1}^{n}A_i$.
    \part For $n>2$, the events $A_i$ are mutually exclusive, ignoring cases with probability zero. Starting at a point $P_{i_0}$, if $A_{i_0}$ is true, then all other $n-1$ points are in the semicircle going clockwise drawn from $P_{i_0}$. Taking any point $P_{i_1}$ of these $n-1$ points, $P_{i_0}$ is in the semicircle going counterclockwise drawn from $P_{i_1}$, so $P_{i_0}$ is not in the semicircle going clockwise drawn from $P_{i_1}$. So, $A_{i_1}$ does not hold. Therefore, the events $A_i$ are mutually exclusive excluding probability zero cases.
    \part For a fixed $P_i$, the probability that all $n-1$ other points are on the semicircle going clockwise from $P_i$ is $\frac{1}{2^{n-1}}$. So, $P(A_i)=\frac{1}{2^{n-1}}$. $P(A)=P\left(\bigcup_{i=1}^{n}A_i\right)$. Since the overlaps have probability zero, $P\left(\bigcup_{i=1}^{n}A_i\right)=\sum_{i=1}^{n}P(A_i)=nP(A_i)=\frac{n}{2^{n-1}}$.
    \newline
\end{parts}

\underline{Problem 5}
\newline
\begin{parts}
    \part $f_{XY}(x,y)\ge0$ for all $x,y$ in its domain. The domain of the pdf is the triangle with vertices at $(0,0),(1,0),(0,1)$. Integrating over this domain with $f_{XY}(x,y)$, we have
    \[\int_0^1\int_0^{1-y}(24xy)dx\,dy\]
    \[=\int_0^1(12(1-y)^2y)dy\]
    Using $u=1-y$ and changing the variable back to $y$, we have
    \[=12\int_0^1(y^2(1-y))dy\]
    \[=12\left(\frac{1}{3}-\frac{1}{4}\right)\]
    \[=12\left(\frac{1}{12}\right)\]
    \[=1\]
    Since this is equal to 1, $f_{XY}(x,y)$ is a joint probability density function.
    \part \[E(X)\]
    \[=\int_0^1\int_0^{1-y}(x)(24xy)dx\,dy\]
    \[=\int_0^1(8(1-y)^3y)dy\]
    Using $u=1-y$ and changing the variable back to $y$, we have
    \[=8\int_0^1(y^3(1-y))dy\]
    \[=8\left(\frac{1}{4}-\frac{1}{5}\right)\]
    \[=8\left(\frac{1}{20}\right)\]
    \[=\frac{2}{5}\]
    \part \[E(Y)\]
    \[=\int_0^1\int_0^{1-y}(y)(24xy)dx\,dy\]
    \[=\int_0^1(12(1-y)^2y^2)dy\]
    \[=12\int_0^1(y^2-2y^3+y^4)dy\]
    \[=12\left(\frac{1}{3}-\frac{2}{4}+\frac{1}{5}\right)\]
    \[=12\left(\frac{1}{30}\right)\]
    \[=\frac{2}{5}\]
    \newline
\end{parts}

\underline{Problem 6}
\newline
\begin{parts}
    \part For $X$ and $Y$ to be independent, $f_X(x)f_Y(y)=f_{XY}(x,y)$ must be true.
    \[f_X(x)\]
    \[=\int_{0}^{1}(x+y)dy\]
    \[=x+\frac{1}{2}\]
    \[f_Y(y)\]
    \[=\int_{0}^{1}(x+y)dx\]
    \[=y+\frac{1}{2}\]
    \[f_X(x)f_Y(y)=\left(x+\frac{1}{2}\right)\left(y+\frac{1}{2}\right)\ne x+y=f_{XY}(x,y)\]
    So, $X$ and $Y$ are not independent.
    \part From the last part, $f_X(x)=x+\frac{1}{2}$, for $0\le x\le1$.
    \part We are looking for $P(X+Y<1)$. The domain of the pdf is the unit square with vertices at $(0,0),(1,0),(1,1),(0,1)$. We are looking at the region where $Y<1-X$, which is the triangle with vertices at $(0,0),(1,0),(0,1)$. Integrating over the region with $f_{XY}(x,y)$, we have
    \[\int_{0}^{1}\int_{0}^{1-y}(x+y)dx\,dy\]
    \[=\int_{0}^{1}\left(\frac{1}{2}(1-y)^2+(1-y)y\right)dy\]
    \[=\frac{1}{6}+\frac{1}{2}-\frac{1}{3}\]
    \[=\frac{1}{3}\]
    So, $P(X+Y<1)=\frac{1}{3}$.
    \newline
\end{parts}

\underline{Problem 7}
\newline
\begin{parts}
    \part \[F_{ABC}(a,b,c)=P(A\le a,B\le b,C\le c)\]
    By independence, we have, for $0<a,b,c<1$,
    \[=P(A\le a)P(B\le b)P(C\le c)\]
    \[=abc\]
    \part For $Ax^2+Bx+C=0$ to have all real roots, $B^2-4AC\ge0$ must be true. So, we are looking for $P(B^2-4AC\ge0)$. This is equivalent to $P(B\ge2\sqrt{AC})$. So, just in the $B$-direction, we are integrating from $2\sqrt{AC}$ to 1. Now, we need to integrate over the correct region in the $AC$-plane. We need to first ensure that $2\sqrt{AC}\le1$. This gives, $AC\le\frac{1}{4}$. When looking at the graph of the region, on $0<A<\frac{1}{4}$, the upper bound of $C$ is 1, while on $\frac{1}{4}<A<1$, the value of $C$ is $\frac{1}{4A}$. So, the region can be split into two integrals. So, on the $AC$-plane, the sum of the integrals is $\int_{0}^{\frac{1}{4}}\int_{0}^{1}dc\,da+\int_{\frac{1}{4}}^1\int_{0}^{\frac{1}{4A}}dc\,da$. After incorporating the integral for the $B$-direction, the final result is $P(B^2-4AC\ge0)=\int_{0}^{\frac{1}{4}}\int_{0}^{1}\int_{2\sqrt{AC}}^{1}db\,dc\,da+\int_{\frac{1}{4}}^1\int_{0}^{\frac{1}{4A}}\int_{2\sqrt{AC}}^{1}db\,dc\,da\approx0.254$.
    \newline
\end{parts}

\underline{Problem 8}
\newline
\begin{parts}
    \part Let $X$ and $Y$ be the amount of time it takes to service A.J.'s and M.J.'s car, respectively. We are looking for $P(Y+t<X)$.
    \[P(Y+t<X)\]
    \[=\int_{t}^{\infty}\int_{0}^{x-t}e^{-(x+y)}dy\,dx\]
    \[=\int_{t}^{\infty}e^{-x}(1-e^{-(x-t)})dx\]
    \[=\int_{t}^{\infty}(e^{-x}-e^{-2x}e^t)dx\]
    \[=e^{-t}-e^t(\frac{1}{2}e^{-2t})\]
    \[=\frac{1}{2}e^{-t}\]
    \part We are looking for $P(X+Y<2)$.
    \[P(X+Y<2)\]
    \[=\int_{0}^{2}\int_{0}^{2-y}e^{-(x+y)}dx\,dy\]
    \[=\int_{0}^{2}e^{-y}(-e^{-(2-y)}+1)dy\]
    \[=\int_{0}^{2}(e^{-y}-e^{-2})dy\]
    \[=-e^{-2}+e^0-2e^{-2}\]
    \[=1-3e^{-2}\]
    \newline
\end{parts}

\end{document}